\documentclass[aspectratio=169]{beamer}
\usetheme{Madrid}
\setbeamertemplate{title page}[default][colsep=-4bp,rounded=true,shadow=false]
\usepackage{amsmath,graphicx,booktabs}
\setbeamertemplate{navigation symbols}{}
\title{Air Quality Prediction}
\subtitle{CSAI Project 19 --- Phase 1 Machine Learning}
\author{Project 19 Team}
\date{}
\newcommand{\resultroot}{../../artifacts/phase1}
\IfFileExists{configs/phase1.toml}{\renewcommand{\resultroot}{artifacts/phase1}}{}
\begin{document}
\begin{frame}\titlepage\end{frame}

\begin{frame}{What we want to predict}
\begin{itemize}
\item Predict PM2.5 one hour ahead, in $\mu$g/m$^3$.
\item After the 10:00 reading, predict the concentration at 11:00.
\item Use pollution and weather through 10:00, plus the known calendar.
\item Compare XGBoost with persistence, daily seasonal naive, and Ridge.
\end{itemize}
\medskip
Our question: how much does this added complexity improve on the latest reading?
\end{frame}

\begin{frame}{The supplied data and our cleaning decisions}
\begin{itemize}
\item UCI Beijing PM2.5: 43,824 hourly rows, January 2010--December 2014.
\item Sorted, unique timestamps; no missing timestamp rows.
\item 2,067 missing PM2.5 readings (4.72\%); weather fields are complete.
\item 214 PM2.5 gaps; the longest lasts 155 hours.
\item Keep missing measurements and valid peaks. Exclude the row ID.
\item Pollution and airport weather come from different locations.
\end{itemize}
\medskip
Missing targets are never filled for scoring. Ridge imputes inputs inside each training fold; XGBoost uses missing-value splits.
\end{frame}

\begin{frame}{Training EDA: the upper tail matters}
\begin{columns}
\column{0.43\textwidth}
\begin{itemize}
\item Only 2010--2012: 24,418 measured concentrations.
\item Missing PM2.5: 7.17\%.
\item Mean 97.82; median 73.00.
\item Range 0--994; training p90 = 218.
\item Retain 797 readings above the IQR fence of 305.50.
\end{itemize}
\column{0.55\textwidth}
\includegraphics[width=\linewidth]{\resultroot/eda/distribution_seasonality.png}
\end{columns}
\end{frame}

\begin{frame}{What recent history and weather tell us}
\centering
\includegraphics[height=0.58\textheight]{\resultroot/eda/correlation_autocorrelation.png}
\par\smallskip
\begin{itemize}
\item PM2.5 autocorrelation: 0.961 at 1 hour; 0.394 at 24 hours; about $-0.006$ at 168 hours.
\item Accumulated wind-speed correlation: $-0.264$. Association does not establish causation.
\end{itemize}
\end{frame}

\begin{frame}{Turning observations into forecast inputs}
\begin{itemize}
\item PM2.5 lags: 0, 1, 2, 3, 6, 12, 24, 48, 72 and 168 hours.
\item Trailing means, standard deviations and observed fractions.
\item Known target hour, weekday and month: sine/cosine; weekend flag.
\item Numeric weather at offsets 0, 1 and 24; current wind direction.
\item Missingness stays visible. All windows end at origin $t$.
\end{itemize}
\medskip
70 columns before wind encoding; 73 after encoding the four observed categories.\par
Lag zero is valid: our label is $y(t+1)$, while the latest input is $y(t)$.
\end{frame}

\begin{frame}{Why these models and references?}
\begin{itemize}
\item Persistence: $\hat y_{t+1|t}=y_t$.
\item Daily seasonal naive: $\hat y_{t+1|t}=y_{t-23}$.
\item Ridge: squared error with an L2 penalty, fixed $\alpha=10$.
\item XGBoost: additive trees with gradient/Hessian updates and complexity penalties.
\item The readings connect to Liang et al. (pollution/weather), Chen and Guestrin (boosting), Hyndman and Athanasopoulos (evaluation), and Hoerl and Kennard (Ridge).
\end{itemize}
\end{frame}

\begin{frame}{How we keep the comparison chronological}
\begin{center}
\begin{tabular}{ll}
\toprule
2010--2012 & Initial training and EDA\\
2013 & Four expanding validation quarters\\
2014 & Four backtest quarters; frozen settings\\
\bottomrule
\end{tabular}
\end{center}
\begin{itemize}
\item Fit once per quarter using labels available by the first origin.
\item Keep the fit fixed within each quarter; update observed histories hourly.
\item Later quarters can use earlier observed quarters.
\item Every method scores the same observed targets and baseline inputs.
\item 2014 coverage: 8,531 / 8,760 hours, or 97.39\%.
\end{itemize}
\end{frame}

\begin{frame}{Parameter selection and feature comparison}
\begin{columns}
\column{0.48\textwidth}
\begin{itemize}
\item Three candidates: validation RMSE 26.772, 26.801 and 26.953.
\item Selected: 300 trees, depth 3, rate 0.05, minimum child weight 5.
\item Row/column sampling 0.8/0.9, L2 = 5, seed 42.
\item Feature groups use these same settings and 2013 dates only.
\end{itemize}
\column{0.50\textwidth}
\includegraphics[width=\linewidth]{\resultroot/figures/validation_ablation.png}
\end{columns}
\medskip
Full-feature XGBoost still trails validation persistence (24.870 RMSE) and Ridge (23.809).
\end{frame}

\begin{frame}{Measured results on 2014}
\centering
\input{\resultroot/results_table.tex}
\par\medskip
MAE and RMSE are in $\mu$g/m$^3$; all rows use 8,531 common hours.
\begin{itemize}
\item XGBoost improves RMSE over persistence by 4.29\%; Ridge by 4.41\%.
\item Ridge is slightly better overall: its RMSE is lower by only 0.025.
\item No statistical significance claim is made for that small difference.
\end{itemize}
\end{frame}

\begin{frame}{The ranking changes across quarters}
\centering
\includegraphics[height=0.69\textheight]{\resultroot/figures/quarterly_rmse.png}
\par\smallskip
Ridge is better in Q1/Q3; XGBoost in Q2/Q4. Both improve RMSE over persistence in each quarter.
\end{frame}

\begin{frame}{High-pollution hours remain difficult}
\centering
\includegraphics[height=0.57\textheight]{\resultroot/figures/pollution_regime_errors.png}
\begin{itemize}
\item At/above training p90 = 218: 865 common hours.
\item XGBoost RMSE 39.959 versus persistence 38.557; peak bias is $-11.912$.
\item Its overall improvement does not establish better peak forecasting.
\end{itemize}
\end{frame}

\begin{frame}{A week of historical backtest forecasts}
\centering
\includegraphics[width=\textwidth]{\resultroot/figures/forecast_last_week.png}
\par\smallskip
The saved deployment model is a later refit. It was not used to produce these backtest predictions.
\end{frame}

\begin{frame}{What we can reproduce, and what remains uncertain}
\begin{itemize}
\item Saved input/source hashes, versions, configuration, selection record and forecast rows.
\item Cleaned hourly export, gap dates, feature dictionary and worked feature preview.
\item Residual lag-1 correlation: XGBoost 0.120; Ridge 0.063.
\item Historical location only; weather publication delays are assumed.
\item No causal weather claim, calibrated uncertainty or daily AQI label.
\item Holdout results have been viewed; further tuning needs fresh evaluation.
\end{itemize}
\end{frame}

\begin{frame}{What we learned}
\begin{itemize}
\item The latest reading is a strong one-hour forecast.
\item Our XGBoost and Ridge models improve pooled RMSE modestly in 2014.
\item Temporal and weather features help this XGBoost setup on validation dates.
\item Ridge remains a competitive, slightly stronger overall comparator.
\item Peak prediction needs further work; the pooled score hides that weakness.
\end{itemize}
\medskip
This completes the classical ML experiment for Phase 1. Deep learning and a hybrid model belong to the later phase.
\end{frame}
\end{document}
